\subsection{Desarrollo del módulo de captura de escritura}

El módulo de captura de escritura tiene como objetivo permitir que el alumno registre evidencia escrita de los ejercicios realizados para su posterior análisis mediante los módulos OCR, NLP y evaluación caligráfica. Este módulo fue desarrollado utilizando tecnologías web del lado del cliente, integrando herramientas de captura manual, carga de imágenes y acceso a cámara desde el navegador.

\subsubsection{Implementación del lienzo digital}

El módulo de lienzo digital fue desarrollado utilizando React y la API Canvas de HTML5, permitiendo al alumno realizar ejercicios de escritura directamente desde el navegador. La implementación incluye herramientas de lápiz y borrador, así como funcionalidades para deshacer acciones, limpiar el área de trabajo y exportar el dibujo realizado.

Para la captura de trazos se implementaron eventos de mouse y dispositivos táctiles, permitiendo compatibilidad tanto en computadoras como en dispositivos móviles. El contexto del canvas se configura dinámicamente dependiendo de la herramienta seleccionada, ajustando propiedades como color, grosor y tipo de operación gráfica.

El siguiente fragmento muestra la lógica principal utilizada para iniciar y realizar el dibujo dentro del lienzo:

\begin{lstlisting}[language=JavaScript, caption={Implementación básica del lienzo digital}]
const iniciarDibujo = (e) => {
    const { x, y } = obtenerCoordenadas(e);

    if (esBorrador) {
        contextRef.current.globalCompositeOperation = 'destination-out';
        contextRef.current.lineWidth = grosurGoma;
    } else {
        contextRef.current.globalCompositeOperation = 'source-over';
        contextRef.current.strokeStyle = 'black';
        contextRef.current.lineWidth = grosurLapiz;
    }

    contextRef.current.beginPath();
    contextRef.current.moveTo(x, y);
    setIsDrawing(true);
};

const dibujar = (e) => {
    if (!isDrawing) return;

    const { x, y } = obtenerCoordenadas(e);

    contextRef.current.lineTo(x, y);
    contextRef.current.stroke();
};
\end{lstlisting}

Adicionalmente, se implementó un mecanismo de historial basado en imágenes serializadas del canvas, permitiendo restaurar estados anteriores mediante la funcionalidad de deshacer.

Una vez finalizado el ejercicio, el contenido del lienzo se convierte en una imagen codificada en Base64 y se envía al backend mediante una solicitud HTTP autenticada.

\begin{lstlisting}[language=JavaScript, caption={Exportación y envío del dibujo al backend}]
const imagenDataUrl = canvasRef.current.toDataURL('image/png');

await fetch('/api/intentos/registrar', {
    method: 'POST',
    headers: {
        'Content-Type': 'application/json',
        'Authorization': `Bearer ${token}`
    },
    body: JSON.stringify({
        id_ejercicio_tutor: idEjercicioTutor,
        imagen_codificada: imagenDataUrl
    })
});
\end{lstlisting}

\subsubsection{Implementación de carga de imagen y cámara}

El módulo de captura de escritura fue desarrollado en React, permitiendo al alumno registrar evidencias mediante la cámara del dispositivo o la carga de imágenes desde el almacenamiento local. Su objetivo es facilitar la adquisición de muestras de escritura para posteriormente ser procesadas por los módulos OCR y NLP.

La implementación valida el tipo y tamaño de los archivos recibidos, aceptando únicamente imágenes con un límite máximo de 10 MB. Una vez seleccionada la imagen, esta es redimensionada y comprimida utilizando un elemento \texttt{canvas}, optimizando así el envío y almacenamiento de la información.

\begin{lstlisting}[language=JavaScript, caption={Procesamiento y validación de imágenes}]
const procesarArchivo = (file) => {
    if (!file.type.startsWith('image/')) return;
    if (file.size > 10 * 1024 * 1024) return;

    const reader = new FileReader();

    reader.onloadend = () => {
        const img = new Image();

        img.onload = () => {
            const canvas = document.createElement('canvas');
            canvas.width = 1200;
            canvas.height = 800;

            canvas.getContext('2d')
                  .drawImage(img, 0, 0, 1200, 800);

            setImagenPreview(
                canvas.toDataURL('image/jpeg', 0.8)
            );
        };

        img.src = reader.result;
    };

    reader.readAsDataURL(file);
};
\end{lstlisting}

Para la captura desde cámara web se utilizó la API \texttt{navigator.mediaDevices.getUserMedia()}, permitiendo acceder al dispositivo de video y obtener fotografías directamente desde la aplicación.

\begin{lstlisting}[language=JavaScript, caption={Inicialización de la cámara web}]
const iniciarCamaraWeb = async () => {
    const mediaStream =
        await navigator.mediaDevices.getUserMedia({
            video: {
                facingMode: 'environment'
            }
        });

    setStream(mediaStream);
    setMostrarCamaraWeb(true);
};
\end{lstlisting}

Asimismo, se incorporaron herramientas básicas de edición como rotación y recorte de imagen, con el propósito de mejorar la calidad de las capturas antes de ser enviadas al servidor.

\subsubsection{Implementación de procesamiento de imagen}

Antes de enviar las evidencias de escritura al servidor, las imágenes pasan por una etapa de preprocesamiento cuyo objetivo es optimizar su calidad y reducir el tamaño de almacenamiento y transmisión. Este proceso mejora la compatibilidad con los módulos OCR y NLP encargados del análisis posterior.

El preprocesamiento implementado contempla:

\begin{itemize}
    \item Validación del tipo y tamaño del archivo.
    \item Redimensionamiento de la imagen para mantener dimensiones uniformes.
    \item Compresión en formato JPEG para reducir el peso del archivo.
    \item Conversión de la imagen a formato Base64 para su envío mediante HTTP.
\end{itemize}

El redimensionamiento y compresión se realizaron utilizando un elemento \texttt{canvas}, sobre el cual se dibuja la imagen original antes de exportarla nuevamente.

\begin{lstlisting}[language=JavaScript, caption={Preprocesamiento de imagen antes del envío}]
const canvas = document.createElement('canvas');

canvas.width = width;
canvas.height = height;

const ctx = canvas.getContext('2d');

ctx.drawImage(img, 0, 0, width, height);

const imagenProcesada =
    canvas.toDataURL('image/jpeg', 0.8);

setImagenPreview(imagenProcesada);
\end{lstlisting}

Finalmente, la imagen procesada es enviada al backend mediante una solicitud HTTP para su almacenamiento y posterior análisis automático.

